\section{Conclusion}
\label{sec:conclusion}

Adding .NET Aspire orchestration to an existing solution is, in practice, a five-step operation: create the AppHost project with its secondary SDK and two NuGet references; write a \texttt{Program.cs} that chains \texttt{AddPostgres}, \texttt{WithDataVolume}, \texttt{WithPgAdmin}, \texttt{WithLifetime(Persistent)}, \texttt{AddDatabase}, and \texttt{AddProject}; create the ServiceDefaults library and call \texttt{AddServiceDefaults()} in the web project; add \texttt{Aspire.Npgsql.EntityFrameworkCore.PostgreSQL} to the web project and register the DbContext; register both new projects in the solution file. The result is a development environment in which Postgres starts automatically, state survives restarts, pgAdmin is available at a dashboard URL, connection strings are injected without manual configuration, and OpenTelemetry signals from every project appear in a single browser dashboard.

The setup has a non-trivial number of coupling points --- the database resource name must match between AppHost and web project, the \texttt{UserSecretsId} must be set for credential persistence, \texttt{WaitFor} must precede the web project launch, and \texttt{ContainerLifetime.Persistent} must be combined with \texttt{WithDataVolume} to achieve durable state. Each of these coupling points is a potential failure mode, and Table~\ref{tab:builder-calls} and Table~\ref{tab:credential-pitfalls} are intended as quick reference for diagnosing the most common ones.

The broader observation is that the friction historically associated with local database setup --- credential management, connection-string plumbing, container lifecycle, health-check ordering --- is now automatable at the application layer in C\#, with no YAML authoring and no manual Docker commands required for the happy path. For teams whose current onboarding process includes a step labelled ``set up local Postgres,'' the Aspire AppHost pattern eliminates that step entirely.
